number of elements in a minimum system of generators of the corresponding ideal; by the preceding proof, we have
@@M0@@
It follows that the kernel of @@M1@@ is locally free of rank zero, and hence that the sequence (1.6.1) is exact.

ii) We may suppose @@M2@@ free over @@M3@@, and @@M4@@ free over @@M5@@; moreover, the hypothesis on the sequence of conormal sheaves allows us to suppose that @@M6@@ is a direct summand of @@M7@@. We can therefore find elements @@M8@@ of @@M9@@ lifting a basis of @@M10@@, and elements @@M11@@ of @@M12@@ such that @@M13@@ lift a basis of @@M14@@, the images of @@M15@@ generating @@M16@@; after further localization, we may suppose that @@M17@@ form a minimum system of generators of @@M18@@, that @@M19@@ form a minimum system of generators of @@M20@@, and that @@M21@@ generate @@M22@@. We can then construct the diagram (1.7.1) and use the spectral sequence (1.7.3).

By (1.7.2), since @@M23@@ is regular, we have @@M24@@ for @@M25@@; we therefore obtain edge-isomorphisms: @@M26@@. Since @@M27@@ is regular, this implies by (1.7.2) that @@M28@@ for @@M29@@, and hence that @@M30@@ is regular, again by (1.7.2).

iii) This last property is proved in EGA IV 19.1.5 (iv), since in the noetherian case the definition in EGA IV coincides with the one we adopt here.

\sgaStableTarget{sga6.e07.prop.1-8}\textbf{Proposition 1.8.} \uline{Let} @@M31@@ \uline{be a scheme}, @@M32@@ \uline{an} @@M33@@-\uline{Module locally free of finite type}, @@M34@@ \uline{a regular homomorphism}, @@M35@@. \uline{Denote by} @@M36@@
